\section{Antecedentes}

\noindent\textbf{Internacional.} La literatura técnica de los fabricantes de PMMA documenta ventanas de temperatura, tiempos de calentamiento por espesor y enfriamiento bajo restricción \cite{plexiglas_forming,throne2008}. Gilormini \emph{et al.} propusieron ecuaciones constitutivas viscoelásticas no lineales para PMMA cerca de $T_g$, útiles para predecir conformado y recuperación elástica \cite{gilormini2010}; Yan \emph{et al.} caracterizaron el \emph{crazing} bajo esfuerzo y su mitigación por recocido \cite{yan2023}. Comercialmente, Shannon Machines ofrece líneas modulares de \SIrange{500}{3000}{\milli\meter} con hilo resistivo de altura ajustable \cite{shannon_line,shannon_bending}, y varios fabricantes asiáticos venden dobladoras automáticas con selección de ángulo y temporización de ablandamiento \cite{mic_auto}.

\noindent\textbf{Nacional y local.} En Colombia hay una tradición consolidada de trabajos de grado sobre máquinas dobladoras, casi siempre para metales: doblado de tubo por compresión, rodillos y rotación \cite{uniandes2003}, y dobladoras hidráulicas de tubería redonda desarrolladas en EAFIT \cite{eafit_hidraulica}. Aportan metodología de diseño mecánico, dimensionamiento de accionamientos y validación dimensional, pero su fenómeno dominante es la plasticidad del metal, no la transición vítrea de un polímero amorfo, por lo que el control térmico queda abierto. EAFIT dispone además de taller de mecanizado, laboratorios de electrónica, instrumentación y caracterización térmica, y del Laboratorio de Mecatrónica, que planteó la necesidad y recibirá el prototipo.
